% ============================================================
%  SECTION 8: DISCUSSION
% ============================================================

\section{Discussion}
\label{sec:discussion}

\subsection{Georgia: A Regressive Prediction Confirmed}

Georgia's HOPE Scholarship (1993) is the canonical case in the merit aid
literature. \textcite{dynarski2000hope} documents that Georgia HOPE
widened the attendance gap between higher- and lower-income students,
with Black students---disproportionately Pell-eligible---showing no
attendance gain. Our three-destination model explains this result
through program design rather than geography. From 1993 to 2000, Georgia
HOPE awards were reduced dollar-for-dollar by Pell Grant amounts,
eliminating the access margin for the lowest-income families (see
Section~\ref{sec:theory_georgia}). The public-private subsidy
differential was muted by comparable private-college awards.
\textcite{cornwell2006enrollment} confirm that four-year private
enrollment in Georgia \emph{rose} by 13.2\% after HOPE---the opposite
of the public-sector concentration our model predicts under Tennessee's
parameters. Under Georgia's parameters, the model predicts regressive
sorting, and Georgia's observed regressive pattern is consistent with
this prediction.

\subsection{West Virginia: Enrollment Contraction Confounds Composition}

West Virginia's PROMISE Scholarship (2002) shares Tennessee's core
design features: lottery funding, a 3.0 GPA threshold, no Pell offset,
and subsidies concentrated at public institutions
\parencite{scottclayton2011money}. If program design alone determined
the direction of compositional sorting, WV should replicate Tennessee's
progressive result. It does not. The Callaway--Sant'Anna ATT for
bottom-quintile share is $-1.21$ pp ($p < 0.001$), indicating a
regressive shift.

The explanation is straightforward: total enrollment at WV public
colleges declined by 155 students per cohort ($p < 0.001$) during the
analysis period, driven by demographic contraction rather than
compositional sorting. Bottom-quintile count fell by 26 per cohort
($p < 0.001$). The modest TWFE share increase for Q1 is a denominator
artifact of this enrollment decline, not progressive sorting.
Balanced-panel estimates restricted to WV colleges with complete data
yield null results. West Virginia is uninformative as a test of the
sorting mechanism---the institutional context prevents the
three-destination mechanism from operating---rather than a contradiction
of the Tennessee finding.

\subsection{Mobility Implications}
\label{sec:mobility}

A natural concern is whether the progressive composition shift came at
the cost of diminished economic returns. If public colleges absorb more
bottom-quintile students but institutional quality declines, the net
welfare effect is ambiguous. Figure~\ref{fig:tn_mobility_trends} offers
a first, descriptive reassurance. The upward mobility rate---the
probability that a student from the bottom income quintile reaches the
top quintile---at Tennessee public colleges tracked below the national
control average in the pre-HOPE period and converged toward it in the
post-period. Despite substantial access expansion, upward mobility at
Tennessee public colleges did not decline.

\begin{figure}[H]
  \centering
  \includegraphics[width=0.85\textwidth]{fig_tn_mobility_trends.pdf}
  \caption{Upward Mobility Rates: Tennessee Public Colleges vs.\
    Control States}
  \label{fig:tn_mobility_trends}
  \vspace{1mm}
  \begin{minipage}{0.85\textwidth}
    \footnotesize
    \textit{Notes:} Mobility rate is the probability that a student
    from the bottom income quintile reaches the top quintile
    (mr\_kq5\_pq1). Group 1 = Tennessee public colleges; Group 0 =
    control-state public colleges. The vertical dashed line marks
    treatment onset. This is a descriptive comparison, not a causal
    estimate. Data: \textcite{chetty2020mobility}.
  \end{minipage}
\end{figure}

We complement this descriptive evidence with an instrumental variables
strategy that provides suggestive evidence on whether composition
shifts are associated with changes in downstream mobility. We
construct a shift-share (Bartik) instrument that predicts each
college's bottom-quintile share using the interaction of its
pre-HOPE income composition with commuting-zone-level trends in the
income distribution of college-going families.

The first stage is strong. The CZ-level Bartik instrument predicts
bottom-quintile share with a Kleibergen--Paap $F$-statistic of 119,
well above conventional weak-instrument thresholds.
Figure~\ref{fig:iv_decomposition} displays the actual and
Bartik-predicted composition paths for Tennessee public colleges. The
two series track closely in the pre-period and diverge after HOPE:
actual bottom-quintile share rises above the demographically predicted
level, with the wedge representing HOPE's policy-attributable
contribution to the progressive shift.

\begin{figure}[H]
  \centering
  \includegraphics[width=0.85\textwidth]{fig_iv_decomposition.pdf}
  \caption{Actual vs.\ Bartik-Predicted Composition: Tennessee Public
    Colleges}
  \label{fig:iv_decomposition}
  \vspace{1mm}
  \begin{minipage}{0.85\textwidth}
    \footnotesize
    \textit{Notes:} Bartik prediction uses pre-HOPE college quintile
    shares $\times$ commuting-zone-level income trends.
    The shaded area between actual and predicted
    composition is the policy-attributable shift. Vertical dashed line
    marks treatment onset ($t^* = 1986$). Data:
    \textcite{chetty2020mobility}.
  \end{minipage}
\end{figure}

Table~\ref{tab:iv_second_stage} reports 2SLS estimates for four child
outcomes, instrumenting bottom-quintile share with the CZ-level
Bartik. Instrumented increases in bottom-quintile representation are
associated with higher upward mobility ($0.114$, SE $= 0.021$,
$p < 0.001$) and lower zero-earnings rates ($-0.225$, SE $= 0.061$,
$p < 0.001$). Mean child earnings rank rises ($0.136$, $p = 0.089$),
and average outcomes among observed bottom-quintile students do not
decline ($0.141$, $p = 0.057$). The sign pattern is coherent across
all four outcomes: no measure suggests that expanded access worsened
downstream returns.

\begin{table}[htbp]\centering
\def\sym#1{\ifmmode^{#1}\else\(^{#1}\)\fi}
\caption{IV Estimates: Effect of Composition on Child Outcomes}
\label{tab:iv_second_stage}
\begin{tabular}{l*{4}{c}}
\toprule
                    &\multicolumn{1}{c}{(1)}&\multicolumn{1}{c}{(2)}&\multicolumn{1}{c}{(3)}&\multicolumn{1}{c}{(4)}\\
                    &\multicolumn{1}{c}{Mobility (Q1$\to$Q5)}&\multicolumn{1}{c}{Child Rank}&\multicolumn{1}{c}{Child Rank $\mid$ Q1}&\multicolumn{1}{c}{Zero Earnings}\\
\midrule
Share of parents in bottom income quintile&       0.114\sym{***}&       0.136         &       0.141\sym{*}  &      -0.225\sym{***}\\
                    &    (0.0208)         &    (0.0799)         &    (0.0742)         &    (0.0606)         \\
\midrule
Observations        &       5,962         &       5,962         &       5,962         &       5,962         \\
KP F-stat           &      118.96         &      118.96         &      118.96         &      118.96         \\
\bottomrule
\multicolumn{5}{l}{\footnotesize Standard errors in parentheses}\\
\multicolumn{5}{l}{\footnotesize 2SLS with college and cohort FEs. SEs clustered by state.}\\
\multicolumn{5}{l}{\footnotesize Endogenous: par\_q1. Instrument: CZ-level Bartik (pre-Q1 $\times$ CZ trend).}\\
\multicolumn{5}{l}{\footnotesize Pooled TN (t*=1986) and MA (t*=1987) public colleges.}\\
\multicolumn{5}{l}{\footnotesize \sym{*} \(p<0.10\), \sym{**} \(p<0.05\), \sym{***} \(p<0.01\)}\\
\end{tabular}
\end{table}


The 2SLS estimates are substantially larger than the corresponding OLS
coefficients, which could reflect negative selection in OLS (colleges
gaining Q1 students through self-selection tend to be lower-quality),
attenuation from measurement error, or the fact that the instrument
identifies a local average treatment effect at the margin of colleges
whose composition shifts because of local demographic change. We do
not take a strong position on which source dominates.

Several limitations warrant caution.  First, the exclusion restriction
requires that CZ income trends affect mobility outcomes only through
college composition. We cannot fully rule out the possibility that
local demographic shifts also directly affect college quality, local
labor-market conditions, or other determinants of student outcomes.
The estimates should therefore be interpreted as evidence on the
consequences of composition shifts induced by local demographic
change, rather than as a fully general estimate of the causal effect
of increasing low-income representation.

Second, pre-trend tests reveal that the Bartik instrument predicts
mobility outcomes in the pre-treatment period as well (joint $F =
20.6$, $p < 0.001$ for upward mobility). This suggests that the
instrument captures pre-existing correlations between local income
trends and college-level mobility, not solely post-treatment
compositional effects. Adding state-by-cohort fixed effects---which
absorb all state-level time-varying confounds---attenuates the
mobility coefficient from $0.114$ to $0.096$ while the first-stage
$F$-statistic remains strong at 85, and the zero-earnings result
is similarly stable ($-0.195$). Mean child rank, however, changes
sign under this more demanding specification, indicating sensitivity
to the control structure.

Third, outcomes defined with respect to the bottom-quintile subgroup
(upward mobility, Q1-conditional earnings rank) are subject to a
compositional interpretation: when bottom-quintile share rises, the
set of enrolled Q1 students may change, and the measured outcome may
reflect this changing composition rather than a treatment effect on a
stable population.

Leave-one-out college diagnostics show that no single institution
drives the results: the mobility coefficient ranges from $0.113$ to
$0.114$ across all 51 leave-one-out iterations. Enrollment-weighted
estimates ($0.127$) are similar to unweighted estimates ($0.114$).

Taking these results together, the IV estimates provide suggestive
evidence that instrumented composition shifts do not reduce downstream
mobility outcomes. The consistency across four conceptually distinct
measures---upward mobility, mean child rank, Q1-conditional rank, and
zero-earnings risk---strengthens this interpretation even though the
instrument does not satisfy a clean exclusion restriction. The most
conservative reading is that colleges whose bottom-quintile
representation rose for instrumented reasons did not experience worse
average outcomes, whether measured by upward mobility, child earnings,
or labor force attachment.

\subsection{Gender Heterogeneity and the Eligibility Channel}
\label{sec:gender_mechanism}

Our three-destination model predicts that compositional sorting operates
through the price channel: the scholarship $S$ reduces the effective
cost of public college for students whose ability exceeds the
eligibility threshold $\tau$. We exploit a sharp implication of this
mechanism. Tennessee HOPE requires a 3.0 high school GPA. Female
students clear this threshold at substantially higher rates than male
students---a well-documented consequence of the persistent gender gap
in high school academic performance
\parencite{goldin2006homecoming}.\footnote{Nationally, the female
  high school GPA distribution is shifted rightward by approximately
  0.2--0.3 grade points. Among 2004 high school graduates,
  approximately 40\% of female students held GPAs above 3.0, compared
  to 30\% of male students \parencite{nces2023digest}.} If HOPE's
composition effect operates through the eligibility-price channel,
the progressive shift should be larger for the gender with a higher
eligibility rate.

We test this prediction using gender-specific parental income measures
from \textcite{chetty2020mobility}. We pool Tennessee and Massachusetts
as treated states (staggered at $t^* = 1986$ and $t^* = 1987$
respectively) against all remaining states and estimate a stacked
interaction model with college-by-gender fixed effects,
cohort-by-gender fixed effects, and state-specific linear gender
trends.\footnote{State-specific linear gender trends absorb
pre-existing differential drift in male versus female composition
across states, which is present in the raw data and produces
misleading pre-trend violations without detrending. After detrending,
the pre-treatment gender interaction is statistically
indistinguishable from zero ($p = 0.913$).} The female-specific
progressive shift is $0.87$ percentage points larger than the male
shift (SE $= 0.0010$, $p < 0.001$). The effect grows over time:
the gender differential is $+0.75$ percentage points in the first
three post-treatment cohorts ($p < 0.001$) and $+0.98$ percentage
points in later cohorts ($p = 0.002$).

The gender differential concentrates in the bottom-quintile share rather
than the mean parental income rank. The mean parental income rank gender gap is $-0.16$
percentage points ($p = 0.130$), statistically indistinguishable from
zero. This pattern is consistent with the eligibility channel operating
primarily at the access margin. Low-income women newly eligible for HOPE
face the sharpest change in their effective choice set: without HOPE,
many would not attend a four-year public institution; with HOPE, the
price drops to zero. Higher-income students attend college regardless
of HOPE, so the gender-eligibility mechanism has little traction in the
middle and upper quintiles. The resulting composition shift is
concentrated in the tail of the income distribution---visible in share
measures but diluted in the mean.

This test has a key identification advantage over the main DiD
specification. Male and female students attend the same colleges in the
same birth cohorts. Any college-level confound---changes in local
economic conditions, institutional quality, or admissions
policy---affects both genders equally and is absorbed by the
college-by-gender fixed effects. Any cohort-level confound---federal
financial aid changes, national labor market shifts---is similarly
absorbed. For a secular trend to produce the observed gender gap, it
would need to differentially shift the income composition of female
versus male students at Tennessee public colleges relative to control
states, in a pattern that begins precisely at the HOPE treatment date.
No plausible alternative satisfies this requirement.
Figure~\ref{fig:gender_event} displays the gender-specific event
studies. Both series track near zero in the pre-period and diverge after
treatment onset, with the female series rising more steeply.

\begin{figure}[H]
  \centering
  \includegraphics[width=0.95\textwidth]{fig_gender_event_study.pdf}
  \caption{Gender Heterogeneity in HOPE's Compositional Effect}
  \label{fig:gender_event}
  \vspace{1mm}
  \begin{minipage}{0.95\textwidth}
    \footnotesize
    \textit{Notes:} Callaway--Sant'Anna doubly-robust event-study
    estimates for bottom-quintile share (Panel~A) and mean parental
    income rank (Panel~B), estimated separately for male and female
    students. The vertical dashed line marks treatment onset
    ($t^* = 1986$). Reference period $k = -1$. Standard errors
    clustered at the state level. Gender-specific outcomes from
    \textcite{chetty2020mobility}, Table~5.
  \end{minipage}
\end{figure}

The gender result provides our most direct test of the
theoretical mechanism. The three-destination model's price channel
predicts a specific pattern---larger progressive shifts for the gender
with higher eligibility rates---and the data confirm it with high
statistical precision. This evidence complements the cross-state
comparisons in the following subsections, which test the model's
predictions along the program-design dimension rather than the
eligibility dimension.

\subsection{Massachusetts Adams Scholarship}

The Adams Scholarship (2004) provides tuition waivers at Massachusetts
public colleges for students scoring in the top 25\% of their district
on the MCAS exam \parencite{cohodes2014merit}. Like Tennessee's HOPE,
the subsidy is concentrated at public institutions---private colleges
receive no benefit---creating a public-private price differential. The
eligibility threshold (top 25\% of district) is moderate, qualifying a
broad share of the income distribution.

The Callaway--Sant'Anna ATT for bottom-quintile share at Massachusetts
public colleges is $+1.38$ percentage points ($p < 0.001$), remarkably
close to Tennessee's $+1.27$ pp. Mean parental income rank declines
by $1.08$ percentage points ($p < 0.001$). The Massachusetts result
provides independent validation of the Tennessee finding: a second
state with a moderate eligibility threshold and a public-sector subsidy
concentration produces a progressive compositional shift of similar
magnitude.

\subsection{South Dakota Opportunity Scholarship}

South Dakota's Opportunity Scholarship requires a minimum ACT score of
24 (approximately the 75th percentile nationally) and provides a
comparatively small award (\$1,000--\$5,000 per year). The high
threshold and small subsidy satisfy the model's conditions for
regressive sorting: only upper-income students are likely to qualify,
and the price differential is too small to induce meaningful sector
switching.

The Callaway--Sant'Anna ATT confirms this prediction: bottom-quintile
share at South Dakota public colleges \emph{declines} by $1.47$
percentage points ($p < 0.001$), and mean parental income rank
\emph{increases} by $3.25$ percentage points ($p < 0.001$). The
effect is strongly regressive---the opposite of Tennessee and
Massachusetts. The quintile decomposition is also inverted: Q1 through
Q3 lose share, Q4 and Q5 gain share.

\subsection{Four-State Comparison}

Figure~\ref{fig:four_state} displays event-study estimates for
bottom-quintile share across all four states. The visual contrast is
striking: Tennessee and Massachusetts trend upward after treatment
(progressive), South Dakota trends sharply downward (regressive), and
West Virginia shows no clear pattern (null, confounded by enrollment
contraction).

\begin{figure}[H]
  \centering
  \includegraphics[width=0.95\textwidth]{fig_four_state_comparison_combined.pdf}
  \caption{Cross-State Evidence: Merit Aid and Institutional Composition}
  \label{fig:four_state}
  \vspace{1mm}
  \begin{minipage}{\textwidth}
    \footnotesize
    \textit{Notes:} Callaway--Sant'Anna event-study estimates for
    bottom-quintile share at public colleges. Tennessee and Massachusetts
    show progressive compositional shifts; South Dakota shows a regressive
    shift; West Virginia shows no significant pattern. Vertical bars
    denote 95\% confidence intervals. Standard errors clustered at the
    state level. Data: \textcite{chetty2020mobility}.
  \end{minipage}
\end{figure}

Table~\ref{tab:design_direction} maps program design parameters to
predicted and observed sorting direction. The three-destination model
correctly predicts the direction of compositional change in all four
cases where program parameters and institutional context are accounted
for.

\begin{table}[H]
\centering
\begin{threeparttable}
\caption{Program Design and Direction of Compositional Sorting}
\label{tab:design_direction}
\begin{tabular}{llccl}
\toprule
State & Program & Threshold & Public Subsidy & Direction \\
\midrule
Tennessee & HOPE & Moderate (3.0 GPA) & Large & Progressive \\
Massachusetts & Adams & Moderate (top 25\%) & Large & Progressive \\
Georgia & HOPE & Moderate (3.0 GPA) & Muted (Pell offset) & Regressive \\
South Dakota & Opp.\ Scholarship & High (ACT 24+) & Small & Regressive \\
West Virginia & PROMISE & Moderate (3.0 GPA) & Large & Null$^\dagger$ \\
\bottomrule
\end{tabular}
\begin{tablenotes}[flushleft]
\footnotesize
\item \textit{Notes:} ``Threshold'' indicates the stringency of
  academic eligibility. ``Public Subsidy'' indicates the size and
  sector-concentration of the scholarship. ``Direction'' reports the
  observed sign of the bottom-quintile share effect. $^\dagger$West
  Virginia's null result reflects enrollment contraction rather than
  a failure of the sorting mechanism
  (Section~\ref{sec:discussion}).
\end{tablenotes}
\end{threeparttable}
\end{table}

\subsection{Generalizability}

The four-state comparison extends the evidence beyond a single treated
state and addresses the concern that Tennessee's result may be
idiosyncratic. Massachusetts independently produces a progressive shift
of similar magnitude under similar program parameters, while South
Dakota's regressive shift confirms the model's prediction that high
thresholds and small subsidies generate the opposite pattern. The
three-destination model organizes all five cases (including Georgia's
documented regressive pattern) through a common set of design
parameters: eligibility threshold location, public-private subsidy
differential, and need-based aid interactions.

The WV non-replication highlights an important scope condition:
program design is necessary but not sufficient. Enrollment contraction
at WV public colleges overwhelmed any compositional dynamics,
suggesting that the model's predictions require a stable or growing
public system with absorptive capacity for the sorting mechanism to
operate.
